%% Week 45 -- Tuesday session (flipped classroom): finalising Project 2
%% Companion to week45slides.tex (Monday lecture on recurrent networks).
%% Two cases only, by design, and no warm-up frames: the first twenty
%% minutes are open questions from Monday (the hidden state, the unrolled
%% graph, BPTT, the Jacobian product and the memory horizon, the LSTM gates,
%% the adding problem), then the rest of the double session is about
%% Project 2 and nothing else: Case 1 = the regression half (parts b-d) --
%% the network against OLS/Ridge/Lasso from Project 1, the hyperparameter
%% heat maps, cross-validated as in week 44; Case 2 = the classification
%% half and the evaluation (parts e-f) -- MNIST, the logistic-regression
%% baseline, the confusion matrix of week 43, honest comparison.  Each case
%% is followed by a "do you have the insight?" frame (answers on \pause).
%% There are NO weekly exercises this week; Project 2 is due Monday
%% November 9 (deadline moved from Friday November 6 on 2026-11-01).  Sources: doc/BookML/Projects/Project2.tex (2026 version,
%% required implementation/analysis lists), the Tuesday decks of weeks
%% 40-44, Sections 2.12, 3.13-3.14, 4.7-4.11, 5.3, 5.6, 5.12, 8.5-8.15 of
%% the book.
%% Build:  pdflatex week45tuesday.tex  (twice)
\documentclass[10pt,aspectratio=169]{beamer}
\usetheme{red_plain}
\usepackage[T1]{fontenc}
\usepackage[utf8]{inputenc}
\usepackage{amsmath,amssymb,bm}
\usepackage{graphicx}
\usepackage{listings}
\usepackage{xcolor}
\usepackage{booktabs}
\graphicspath{{../BookFigures/chapter11_recurrent_networks/}{./}}

%% ---- code listings
\definecolor{codebg}{rgb}{0.97,0.97,0.97}
\lstset{language=Python,basicstyle=\ttfamily\scriptsize,backgroundcolor=\color{codebg},
  keywordstyle=\color{blue!70!black}\bfseries,commentstyle=\color{green!45!black},
  stringstyle=\color{red!60!black},showstringspaces=false,frame=single,
  framerule=0pt,xleftmargin=4pt,xrightmargin=4pt,breaklines=true,columns=fullflexible}

%% ---- discussion frames (kept for future weeks, unused in this deck)
\newcounter{insight}
\newenvironment{insight}[1][]{%
  \stepcounter{insight}%
  \begin{frame}{Warm-up \arabic{insight}: #1}%
  \setbeamercolor{block title}{bg=structure.fg,fg=white}%
  \setbeamercolor{block body}{bg=structure.fg!8}%
}{\end{frame}}
\newcommand{\question}[1]{\begin{block}{Question}#1\end{block}}
\newcommand{\answer}[1]{\pause\begin{block}{Answer}#1\end{block}}
\newcommand{\hint}[1]{\pause{\small\emph{Hint: }#1}}

%% ---- case frames
\newcounter{case}
\newenvironment{case}[1][]{%
  \stepcounter{case}%
  \begin{frame}[fragile,environment=case]{Case \arabic{case}: #1}%
}{\end{frame}}

\setbeamertemplate{footline}[frame number]
\setbeamertemplate{blocks}[rounded][shadow=false]
\setbeamercovered{invisible}

\title{Week 45, Tuesday Session}
\subtitle{Flipped classroom: your questions, then Project 2 only --
  the regression half against Project 1, the classification half against
  logistic regression, and how to write it up}
\author{Morten Hjorth-Jensen}
\institute{Department of Physics and Center for Computing in Science Education,\\
  University of Oslo, Norway}
\date{}

\begin{document}

\begin{frame}
\titlepage
\end{frame}

\begin{frame}{How today works}
\small
\begin{block}{The format (as every Tuesday)}
Mondays are lectures; Tuesdays belong to your questions and to concrete
cases.  Today the cases are \emph{your own results}: the project you hand
in on Monday November 9.
\end{block}
\begin{enumerate}
  \item \textbf{First $\sim$20 min:} your open questions from Monday's
        lecture -- the hidden state and the unrolled graph, backpropagation
        through time, the product of Jacobians and the memory horizon, the
        LSTM gates and the forget-gate bias, the adding problem -- discussed
        in plenum.
  \item \textbf{Rest of the double session:} \emph{two} cases, worked in
        your project groups with your Project 2 code, figures and draft
        open.  Case 1: the regression half (parts b--d) -- your network
        against your own OLS, Ridge and Lasso from Project 1, and the
        hyperparameter maps.  Case 2: the classification half and the
        evaluation (parts e--f) -- MNIST against the logistic-regression
        baseline, the confusion matrix, and the critical comparison.  Each
        case ends with the questions a grader -- or an examiner -- will ask
        you.
  \item \textbf{Last 5 min:} wrap-up, how to hand in, and what Monday builds
        on.
\end{enumerate}
\begin{block}{No weekly exercises this week}
The week is for finishing the project.  \textbf{Project 2 is due Monday
November 9.}
\end{block}
\end{frame}

\begin{frame}{Picking up from Monday -- and where you should be with Project 2}
\footnotesize
\begin{columns}[T]
\begin{column}{0.44\textwidth}
\begin{block}{Monday in one sentence each}
\scriptsize
A CNN shares parameters in space, a recurrent network in time: it carries
a state forward and reuses $\bm{U},\bm{W},\bm{V}$ at every step -- a
learned time-stepper, Eq.~(11.13); in Keras and PyTorch one layer fed a
\texttt{(samples, timesteps, features)} tensor, in the same
\texttt{Sequential}/training-loop pattern as the CNN.  Unrolled, it is a
deep network with tied weights, and \texttt{loss.backward()} is BPTT,
Eqs.~(11.15)--(11.18).  The gradient across a lag is a product of Jacobians
(Theorem 11.1): $10^{-18}$ over $59$ steps at $\rho(\bm{W})=0.5$.  Clipping
fixes explosion only.  The simple RNN learned the oscillator's rule and
fails on the adding problem from $T=30$; gated cells next Monday.
\end{block}
\begin{block}{Questions we expect}
\scriptsize
Why is the second term of Eq.~(11.16) absent in the sequence-to-label case?
Why $\tanh$ and not ReLU?  Why did the network trained on one trajectory
predict a different one?  Why does PyTorch print $673$ parameters where
Keras prints $649$?  Ask.
\end{block}
\end{column}
\begin{column}{0.54\textwidth}
\begin{block}{Project 2: what should be done by now}
\scriptsize
\setlength{\tabcolsep}{3pt}
\begin{tabular}{@{}lll@{}}
\toprule
Part & Done in & Status\\
\midrule
a) costs, activations, derivatives & wk 40--41 & written up\\
b) own network on Runge, vs OLS & wk 41--42 & figures final\\
c) activations, depth, width & wk 42 & heat maps\\
d) $L_2$/$L_1$ vs Ridge/Lasso & wk 42 & table\\
e) MNIST, baseline, confusion matrix & wk 43 & \textbf{Case 2 today}\\
cross-validated hyperparameters & wk 44 & \textbf{Case 1 today}\\
f) evaluation, abstract, conclusion & now & drafted\\
\bottomrule
\end{tabular}
\end{block}
\begin{block}{What remains}
\scriptsize
The \emph{analysis} -- the six items of ``Required analysis'' in the project
text -- and the text around it; a clean re-run from a fresh kernel so that
the numbers in the report are the numbers the code produces; the LLM
appendix; the GitHub link.  Today is about the analysis.
\end{block}
\end{column}
\end{columns}
\end{frame}

\begin{frame}{The analysis is where the points are: what a figure must do}
\footnotesize
\vspace{-8pt}
\begin{columns}[T]
\begin{column}{0.5\textwidth}
\begin{block}{A figure states a finding}
\begin{itemize}
  \item One claim per figure, in the caption, with the number in it:
        ``The network reaches a test MSE of $0.0041\pm0.0006$ at
        $\gamma=10^{-2}$, $\lambda=10^{-4}$, against $0.0039$ for OLS at
        degree 8.''
  \item Logarithmic axes where the quantity spans decades ($\gamma$,
        $\lambda$, MSE).
  \item Heat maps for two hyperparameters at a time (project text, analysis
        item 2): colour = cross-validated score, the chosen cell marked,
        the range wide enough to show \emph{bad} performance on both sides.
  \item Error bars from the folds (week 44) wherever a single number is
        compared with another single number.
\end{itemize}
\end{block}
\end{column}
\begin{column}{0.48\textwidth}
\begin{block}{A table compares on equal terms}
\begin{itemize}
  \item Same data, same split, same scaling, same seed for every row.
  \item Hyperparameters chosen on training data only (validation split or
        $k$-fold, Section 3.14); the test number computed \emph{once} and
        choosing nothing.
  \item Parameter counts next to the scores: a $1$--$50$--$1$ network has
        $151$ parameters, a degree-8 polynomial $9$.
  \item Say what ``best'' means: lowest CV MSE, or simplest within one
        standard error (Section 3.14)?
\end{itemize}
\end{block}
\medskip
\textbf{The sentence after the figure:} its ``so what'' -- what moved, why
(which section of the notes), and what you decided because of it.
\end{column}
\end{columns}
\end{frame}

%% ====================================================================
%% Case 1: the regression half
%% ====================================================================

\begin{case}[the regression half (parts b--d): your network against Project 1]
\footnotesize
\textbf{Goal:} leave with the two figures and the one table that carry the
regression argument of your report, each with a caption that states a
finding, and with the paragraph that answers analysis items 1--3 drafted.
\begin{enumerate}
  \item \textbf{Assemble the evidence} (your own code, the Runge data of
        Project 1): (A) training and test MSE against epochs for the
        $1$--$50$--$1$ and $1$--$100$--$100$--$1$ networks with plain GD,
        RMSprop and Adam at the learning rates you scanned (part b); (B) the
        heat map of CV MSE over $(\gamma,\lambda)$ for your best
        architecture, $L_2$ and $L_1$ separately (part d); (C) the heat map
        over (hidden nodes, activation) or (layers, nodes) (part c).
  \item \textbf{The table:} OLS, Ridge and Lasso at their Project-1 best
        (degree, $\lambda$) against the network at its cross-validated best,
        with CV MSE $\pm$ standard error and the \emph{test} MSE computed
        once -- plus the parameter count of each row.
  \item \textbf{Cross-examine:} does the network beat OLS on a smooth
        one-dimensional function?  (It often does not, and that is a result,
        not a failure.)  Does $L_2$ on the weights behave like Ridge -- a
        shallower, later minimum -- and $L_1$ like Lasso?  Where does the
        sigmoid network stall and the ReLU one not, and what does Section
        8.9 say about it?
  \item \textbf{Draft the paragraph:} which model, why, with the
        bias-variance language of Chapter 2 and the parameter counts.
\end{enumerate}
\end{case}

\begin{frame}{Case 1: what you should find -- and what it looks like when it is wrong}
\footnotesize
\begin{columns}[T]
\begin{column}{0.5\textwidth}
\begin{block}{The picture that should emerge (qualitatively)}
\scriptsize
(A) Plain GD at a fixed $\gamma$ is slow and either rattles or creeps;
Adam reaches the plateau in a fraction of the epochs and is insensitive to
$\gamma$ over a decade (week 38); RMSprop sits between.  The $1$--$50$--$1$
ReLU network reaches a test MSE comparable with the best OLS polynomial of
Project 1 -- often slightly worse, with fifteen times the parameters.  The
deeper network is not better on this problem and overfits sooner.  (B) The
$L_2$ heat map has a valley: too small a $\lambda$ changes nothing, too large
one underfits (the Ridge shrinkage of Eq.~3.48 in a different disguise);
the $L_1$ valley is narrower and the network loses nodes as $\lambda$ grows.
(C) Sigmoid hidden layers train slowest and stall first as the depth grows
(Section 8.9: $\sigma'\le1/4$ per layer); ReLU and leaky ReLU are close;
width beyond $\sim50$ buys little on $n=100$ points.
\end{block}
\end{column}
\begin{column}{0.48\textwidth}
\begin{block}{Symptoms, and what they mean}
\scriptsize
\begin{itemize}
  \item Test MSE far below $\sigma^{2}=0.01$: you scaled with the whole
        data set before splitting, or evaluated on training data
        (Section 3.13).
  \item Network MSE constant in $\gamma$ over four decades: Adam is doing its
        job -- say so -- or the run never left the initialisation; check the
        training curve.
  \item ReLU network at exactly the mean of $y$: dead units from a large
        $\gamma$ or a bad initialisation; use He initialisation (Eq.~8.74)
        and halve $\gamma$.
  \item $L_1$ and $L_2$ results identical: the penalty is in the cost but
        not in the gradient, or $\lambda$ never reached the scale of the
        weights.
  \item Different seed, different ``best'' architecture: the minimum is
        flat; report the region, the error bar and the simplest model inside
        it.
  \item Network beats OLS by a factor ten: check that OLS was given its
        Project-1 best degree and the same scaled data.
\end{itemize}
\end{block}
\end{column}
\end{columns}
\end{frame}

\begin{frame}{Case 1: do you have the insight?}
\scriptsize
\begin{enumerate}
  \item Your $1$--$50$--$1$ network has a cross-validated MSE of
        $0.0112\pm0.0009$; OLS at degree $8$ has $0.0105\pm0.0008$.  Which
        model do you recommend for this problem, and what do you write in the
        conclusion about neural networks for regression?
  \item Adam at $\gamma=10^{-3}$ and at $\gamma=10^{-2}$ give the same final
        MSE; plain gradient descent needs $\gamma=0.5$ and $5000$ epochs to
        get there.  Explain both observations in one sentence each.
  \item With $L_2$ regularisation the heat map's valley sits at
        $\lambda\approx10^{-4}$; Ridge in Project 1 wanted
        $\lambda\approx10^{-2}$.  Are the two $\lambda$'s comparable at all?
  \item The sigmoid network with three hidden layers stays at the mean of
        $y$ for $2000$ epochs and then suddenly learns.  What happened, and
        which figure of the lecture notes would you cite?
\end{enumerate}
\pause
\begin{block}{What we hope you say}
\scriptsize
(1) The two are within one standard error of each other; the simpler model
-- nine parameters, a closed-form solution, no learning rate -- wins on
every count except flexibility.  The honest conclusion is that on a smooth
one-dimensional function with $100$ points a polynomial is a hard
benchmark, and that the network's advantage is that it did not need the
basis to be chosen (Chapter 8, the ``learned features'' insight of the
project text).  (2) Adam divides each component of the gradient by a running
scale, so the step is of order $\gamma$ in every direction and the result is
insensitive to $\gamma$ over a decade (Section 4.11); plain GD must use a
single $\gamma$ below $2/\lambda_{\max}$ for the stiffest direction and then
crawls along the flat ones (Section 4.5).  (3) No: Ridge penalises the
polynomial coefficients, $L_2$ penalises weights and biases of a network
whose scale is set by the initialisation, and the data are scaled
differently; what \emph{is} comparable is the \emph{shape} of the
MSE-against-$\lambda$ curve and the qualitative effect -- a later, shallower
minimum.  Compare mechanisms, not numbers.  (4) Vanishing gradients: three
sigmoid layers multiply the error by $\sigma'\le1/4$ three times, so the
first layer barely moves until the weights have grown enough to carry a
gradient (Section 8.9, Figure 8.14); ReLU or He initialisation removes the
plateau.  Cite the figure and your own training curve side by side.
\end{block}
\end{frame}

%% ====================================================================
%% Case 2: the classification half and the evaluation
%% ====================================================================

\begin{case}[the classification half and the evaluation (parts e--f)]
\footnotesize
\textbf{Goal:} the MNIST table, the confusion matrix with its reading, the
baseline comparison, and the critical paragraph of your conclusion drafted.
\begin{enumerate}
  \item \textbf{Assemble the evidence} (full MNIST, the $60\,000/10\,000$
        split as intended): (A) test accuracy of logistic regression (your
        network with zero hidden layers, softmax + cross entropy from part a)
        -- the required baseline; (B) test accuracy of your best network, its
        hyperparameters chosen on training data only (validation split or
        $k$-fold on a subset); (C) the $10\times10$ confusion matrix of (B)
        on the test set, counts and row-normalised (week 43).
  \item \textbf{The table:} baseline, your network, and -- if you did it --
        Scikit-Learn's \texttt{LogisticRegression} or \texttt{MLPClassifier}
        on the same split, with parameter counts and training time.
  \item \textbf{Read the confusion matrix:} which digits are confused with
        which, and in which direction?  Per-class precision and recall from
        its rows and columns (Section 5.12); which class would you flag to a
        user?
  \item \textbf{Cross-examine the baseline:} \emph{does} a hidden layer
        help, by how many points, and is that difference larger than the
        run-to-run spread (two or three seeds)?  Analysis item 4 asks you to
        say \emph{why}.
  \item \textbf{Draft part f):} regression and classification side by side
        -- where the network earned its parameters and where it did not; the
        gradient methods and learning rates (analysis item 6) in one
        paragraph with the evidence of Case 1.
\end{enumerate}
\end{case}

\begin{frame}{Case 2: what you should find}
\footnotesize
\begin{columns}[T]
\begin{column}{0.5\textwidth}
\begin{block}{The numbers that should emerge}
\scriptsize
Logistic regression on $784$ scaled pixels: about $92\%$ test accuracy
(the baseline named in the project text).  One hidden layer of $100$ ReLU
nodes with softmax and cross entropy, Adam, a few epochs: $97$--$98\%$ --
five to six points above the baseline, and the gap is far larger than the
spread over seeds ($\pm0.2$ points).  Two hidden layers: a few tenths of a
point more at most.  For reference, a CNN reaches $99\%$ (week 44) and an
LSTM reading rows reaches $97\%$ (Monday).  The confusion matrix is
dominated by its diagonal; the off-diagonal mass sits in a few pairs --
$4\leftrightarrow9$, $3\leftrightarrow5$, $7\leftrightarrow9$,
$3\leftrightarrow8$ -- and is \emph{not} symmetric: more $9$s are read as
$4$ than the reverse, say, and that direction is the sentence to write.
Precision and recall per class are within a few points of each other except
for the confused pairs.
\end{block}
\end{column}
\begin{column}{0.48\textwidth}
\begin{block}{Symptoms, and what they mean}
\scriptsize
\begin{itemize}
  \item Accuracy stuck at $11.35\%$: the network predicts one class -- a
        dead ReLU layer or a $\gamma$ too large for the cross entropy; check
        that $\delta^L=\bm{a}-\bm{y}$ is what your code computes.
  \item Training accuracy $100\%$, test $96\%$: overfitting; early stopping
        on a validation split, or $L_2$, or fewer epochs.
  \item Baseline at $98\%$: you trained logistic regression on the test set,
        or the ``baseline'' has a hidden layer.
  \item Confusion matrix with the classes in the wrong order, or rows
        summing to $10\,000$: check \texttt{labels=} and normalise by row for
        recall, by column for precision.
  \item Test accuracy quoted to four decimals: with $10\,000$ test images the
        binomial standard error at $97\%$ is $0.17$ points -- two decimals,
        and say so.
  \item Hyperparameters tuned by looking at the test accuracy: the test
        number then chooses something and is no longer a test number.
        Re-split.
\end{itemize}
\end{block}
\end{column}
\end{columns}
\end{frame}

\begin{frame}{Case 2: do you have the insight?}
\scriptsize
\begin{enumerate}
  \item Your network reaches $97.6\%$ and logistic regression $92.1\%$.  Why
        does a hidden layer help \emph{here} when it did not help in the
        regression half?
  \item Your confusion matrix shows $4$ read as $9$ in $31$ cases and $9$
        read as $4$ in $12$.  Write the sentence for the report, and say
        which of precision and recall of class $4$ is affected by each
        number.
  \item A fellow student reports $98.9\%$ with a dense network trained for
        $200$ epochs on the full training set, hyperparameters chosen by
        watching the test accuracy.  What is wrong, and how large is the
        problem?
  \item Analysis item 6 asks about gradient methods and learning rates.  You
        have the Runge curves of Case 1 and three MNIST runs with Adam at
        $\gamma=10^{-3}$.  Is that enough evidence, and what would you add
        with one more hour of compute?
\end{enumerate}
\pause
\begin{block}{What we hope you say}
\scriptsize
(1) Logistic regression is a linear classifier in pixel space: it draws ten
hyperplanes through $784$ dimensions, and digits are not linearly
separable there -- a $4$ and a $9$ differ by a stroke that no single
weighted sum of pixels isolates.  The hidden layer learns features (strokes,
loops) in which the classes \emph{are} nearly separable, which is the
``learned features'' insight of Chapter 8.  The Runge function, by contrast,
is a smooth function of one variable that a degree-8 polynomial already
represents to the noise level; there was no feature left to learn.  (2)
``Of the $982$ fours, $31$ were classified as nines, while $12$ of the
$1009$ nines were classified as fours'' -- the $31$ lower the \emph{recall}
of class $4$ (row), the $12$ lower its \emph{precision} (column).  (3) The
test set was used to choose hyperparameters, so $98.9\%$ is an optimistic
estimate -- by how much nobody can say from these numbers, which is the
problem: the one number that was supposed to be a prediction of performance
on new data no longer is.  Re-split, tune on a validation set, report the
test number once.  (4) Not quite: the Runge curves show the \emph{mechanism}
(Adam's insensitivity to $\gamma$, plain GD's stability limit), but item 6
asks about training \emph{the} network, which includes MNIST.  One hour: a
$\gamma$ scan over three decades for Adam and RMSprop and two batch sizes on
a $10\,000$-image subset, test accuracy against epoch -- the week-38 figure
on the project's own data.
\end{block}
\end{frame}

\begin{frame}{The last checks before Monday November 9}
\scriptsize
\begin{columns}[T]
\begin{column}{0.5\textwidth}
\begin{block}{Reproducibility}
\begin{itemize}
  \item Restart the kernel, run everything top to bottom, and compare every
        number in the report with the output.  Seeds fixed and stated.
  \item Every figure saved as PDF by the code (\texttt{plt.savefig}), named
        after what it shows, and produced by a script in the repository.
  \item Your tests, quoted: the gradient check of your backpropagation
        against automatic differentiation or central differences (as we did
        on Monday for the recurrent network); a linear network with no hidden
        layer reproducing OLS; agreement with Scikit-Learn or PyTorch on one
        configuration.
  \item Upload the notebook executed top to bottom, so that a reader sees
        the outputs.
\end{itemize}
\end{block}
\end{column}
\begin{column}{0.48\textwidth}
\begin{block}{Hand-in (Canvas, Monday November 9)}
\begin{itemize}
  \item One PDF, the report; the code as a link to your GitHub/GitLab
        repository with a folder of selected results and a README that says
        how to run it.
  \item Cite the data set (LeCun, Cortes and Burges, \emph{The MNIST database
        of handwritten digits}) and the software you used, as the project
        text asks.
  \item The LLM appendix after the conclusions -- missing means $-5$ points
        -- consistent with your docstrings.
  \item References verified against the sources, in the course format.
  \item Group of two to three: all names on the report, one hand-in.
\end{itemize}
\end{block}
\end{column}
\end{columns}
\end{frame}

\begin{frame}{Wrap-up, the deadline, and Monday}
\begin{itemize}
  \item \textbf{No weekly exercises.}  \textbf{Project 2 is due Monday
        November 9} -- PDF on Canvas, code on GitHub, LLM appendix included.
  \item \textbf{If you want Monday's material in the project:} you do not
        have to.  A group that treats the Runge regression as one-step
        forecasting along $x$ and compares a recurrent network with its dense
        network, cross-validated the same way, is welcome to -- as an
        addition, not a replacement.
  \item \textbf{Next Monday (week 46):} decision trees, ensemble methods and
        random forests (Chapter 7), then boosting -- a non-linear model built
        without gradient descent, which will matter for Project 3.
\end{itemize}
\end{frame}

\end{document}
